\documentclass[11pt]{article}
\usepackage{rabbithole}

\phasenumber{4}
\phasetitle{Undo}
\phasesubtitle{Every way to take it back, and the one way you cannot.}

\hypersetup{pdftitle={Rabbit Holes, Phase 4: Undo},
            pdfauthor={Accelerate, Manipal University Jaipur},
            pdfsubject={git and GitHub handbook}}

\begin{document}
\maketitlepage

\section{The premise}

Most people are frightened of git for one reason. They believe that if they type
the wrong thing, work disappears.

This is almost never true, and the gap between ``almost never'' and what people
believe is where all the fear lives. Git is built to hoard. It keeps things you
deleted, commits you abandoned, branches you removed, and the position of your
\co{HEAD} pointer going back ninety days. Undoing something in git is usually not
recovery at all, it is just moving a pointer back to something that never actually
went anywhere.

This phase covers every undo git has. By the end you should be willing to try
things, which matters more than any single command here.

\begin{why}
Confidence is a real engineering skill. A developer who is afraid of their tools
works slowly, avoids branches, commits huge unreviewable lumps, and never
experiments. Learning the undos first is what buys you permission to be careless
later.
\end{why}

\section{The three trees}

Every undo command in git is answering the same question: \emph{which of the three
trees am I changing?}

\begin{description}
  \item[The working directory] The files as they exist on disk right now. What
    your editor shows you.
  \item[The index, also called the staging area] What will go into the next commit.
    \co{git add} copies from the working directory into here.
  \item[HEAD] The commit you currently have checked out. The last recorded state.
\end{description}

When you run \co{git status} and it lists things in two groups, those groups are
``working directory differs from index'' and ``index differs from HEAD''. That is
the whole output.

Hold on to this. Every command below is a different combination of which trees get
touched, and if you know the three trees you can predict what any of them does
without memorising anything.

\section{Undoing work you have not committed}

\subsection{Throw away changes to a file}

You edited a file, you made it worse, you want the version from the last commit.

\begin{shell}
git restore sonnet.md
\end{shell}

The file on disk is replaced with the one in \co{HEAD}.

\begin{trap}
This is the single most dangerous command in this document. The edits you made
were never committed, so git never had a copy, so there is nothing to recover.
This is one of the few genuine ways to lose work in git.

Before running it on something you might want, either commit it or stash it.
Committing on a scratch branch costs nothing and is always reversible.
\end{trap}

\subsection{Unstage a file}

You ran \co{git add} on something you did not mean to include.

\begin{shell}
git restore --staged secrets.env
\end{shell}

This copies the file from \co{HEAD} back into the index. Your edits on disk are
untouched, the file is simply no longer queued for the next commit.

\begin{why}
\co{git restore} was added in Git 2.23, in 2019, together with \co{git switch}.
Before that both jobs were done by \co{git checkout}, which also changed branches,
also restored files, and also created branches. One command, three unrelated jobs,
and the flags that separated them were not memorable.

Older tutorials will tell you \co{git checkout -- file} and \co{git reset HEAD file}.
Both still work. \co{restore} and \co{switch} exist because the maintainers agreed
the old interface was the problem, not you.
\end{why}

\subsection{Delete files git has never seen}

\co{git restore} only knows about tracked files. Build output, stray downloads, and
half-finished experiments are untracked, so it ignores them.

\begin{shell}
git clean -n
git clean -fd
\end{shell}

\co{-n} is a dry run and prints what would be deleted. \co{-f} actually does it,
\co{-d} includes directories.

\begin{trap}
Always run \co{git clean -n} first. Always. This command deletes files git has no
record of, which means there is no version of them anywhere. It is the other
genuine way to lose work.
\end{trap}

\section{Undoing the commit you just made}

\subsection{You want to change the message}

\begin{shell}
git commit --amend -m "A message that is actually useful"
\end{shell}

\subsection{You forgot to include a file}

\begin{shell}
git add the-file-you-forgot.md
git commit --amend --no-edit
\end{shell}

\co{--no-edit} keeps the existing message.

Here is the thing to understand about \co{--amend}: it does not edit anything.
Commits in git are immutable. What actually happens is that git builds a brand new
commit from the current index, gives it the same parent as the old one, and moves
your branch to point at the new commit. The old commit is still sitting in the
object database, unreferenced, and the reflog still knows where it is.

\begin{trap}
Amending rewrites history. If you have already pushed the commit, your local branch
and the remote branch now disagree, and an ordinary push will be rejected. You will
need \co{git push --force-with-lease}, and everyone who already pulled the old
commit will have a bad afternoon.

The rule: amend freely before you push, think hard afterwards.
\end{trap}

\section{git reset, and its three modes}

This is the command people find genuinely confusing, and the confusion is entirely
because it has three modes that do three different amounts of work. The three trees
make it simple.

\co{git reset} moves your branch pointer to a different commit. The mode decides
what happens to the index and the working directory.

\vspace{0.5em}
\begin{tabularx}{\linewidth}{@{}L{2.6cm} X@{}}
\toprule
\rhtablehead{Mode} & \rhtablehead{What it changes} \\
\midrule
\co{--soft} & Moves \co{HEAD} only. Index and working directory untouched. Your
changes are still staged, ready to be recommitted. \\
\co{--mixed} & Moves \co{HEAD} and resets the index. Your changes survive on disk
but are unstaged. \emph{This is the default.} \\
\co{--hard} & Moves \co{HEAD}, resets the index, and overwrites the working
directory. Your changes are gone from all three trees. \\
\bottomrule
\end{tabularx}
\vspace{0.5em}

Read that table as a ladder. Each mode does everything the one above it does, plus
one more tree.

\subsection{Undo the last commit but keep the work}

The one you will use most often. You committed too early, or committed two unrelated
things together.

\begin{shell}
git reset --soft HEAD~1
\end{shell}

The commit is gone, the changes are staged exactly as they were. Fix the grouping,
commit again properly.

\co{HEAD\textasciitilde 1} means ``one commit before HEAD''. \co{HEAD\textasciitilde 3}
means three before. You will see \co{HEAD\^{}} used the same way for the single step
case.

\subsection{Undo the last commit and unstage}

\begin{shell}
git reset HEAD~1
\end{shell}

Same as above but the changes come back unstaged. Useful when you want to re-add
selectively with \co{git add -p}.

\subsection{Destroy the last commit entirely}

\begin{shell}
git reset --hard HEAD~1
\end{shell}

\begin{trap}
\co{--hard} is the flag with a reputation, and it half deserves it. The
\emph{commit} is recoverable through the reflog, covered in a moment. But any
uncommitted changes sitting in your working directory at the time are overwritten,
and those were never in git, so those really are gone.

Committed work: recoverable. Uncommitted work: not. That distinction is the entire
danger and it is worth reading twice.
\end{trap}

\begin{tryit}
Do this in a scratch repo right now, because reading about it is not the same as
watching it.

\begin{shell}
mkdir reset-demo && cd reset-demo && git init
echo one > f.txt && git add . && git commit -m "one"
echo two >> f.txt && git add . && git commit -m "two"

git reset --soft HEAD~1 && git status
git reset --hard HEAD && cat f.txt
git reflog
\end{shell}

Watch \co{git status} after the soft reset, and notice the change is still staged.
Then look at the reflog and find the commit you just destroyed.
\end{tryit}

\section{Undoing something that is already public}

Everything above rewrites history, which is fine on your own machine and rude on a
shared branch. When a bad commit is already on \co{main} and other people have
pulled it, you do not remove it. You add a new commit that reverses it.

\begin{shell}
git revert a1b2c3d
\end{shell}

This creates a fresh commit whose content is the exact inverse of \co{a1b2c3d}.
History gains a commit rather than losing one. Nobody else's clone breaks, because
nothing they already have has changed.

\begin{why}
Rewriting shared history is antisocial for a mechanical reason, not a stylistic one.
Your teammate's branch is built on commits that, after your rewrite, no longer
exist. Their next pull produces a mess that they have to untangle and that they did
not cause.

The standard formulation: \textbf{rewrite your own unpushed history freely, never
rewrite history other people are building on.} Reset for private work, revert for
public work.
\end{why}

To revert a merge commit you have to tell git which side to keep, because a merge
has two parents:

\begin{shell}
git revert -m 1 <merge-sha>
\end{shell}

\co{-m 1} means ``keep the first parent'', which is the branch you merged into,
normally \co{main}.

\section{The reflog}

This is the safety net under every other command in this document, and it is the
single most valuable thing in this phase.

Git records every position \co{HEAD} has occupied. Every commit, checkout, merge,
rebase, reset, and amend appends a line. Not just commits on branches, but every
intermediate state, including ones you have since abandoned.

\begin{shell}
git reflog
\end{shell}

\begin{terminal}
a1b2c3d HEAD@{0}: reset: moving to HEAD~1
9f8e7d6 HEAD@{1}: commit: two
5c4b3a2 HEAD@{2}: commit: one
\end{terminal}

\co{HEAD@\{1\}} is where \co{HEAD} was one move ago. To go back:

\begin{shell}
git reset --hard HEAD@{1}
\end{shell}

Or, better, put the recovered commit somewhere safe rather than moving your branch:

\begin{shell}
git switch -c rescued 9f8e7d6
\end{shell}

\subsection{Recovering a deleted branch}

You deleted a branch you still needed.

\begin{shell}
git reflog
git switch -c the-branch-i-deleted <sha>
\end{shell}

The branch name was only ever a label pointing at a commit. Deleting it removed the
label. The commit was untouched, and the reflog remembers the number.

\subsection{Recovering from a rebase that went wrong}

A rebase leaves a trail. Find the commit from just before you started:

\begin{shell}
git reflog
git reset --hard HEAD@{5}
\end{shell}

Git also stores the pre-rebase position in \co{ORIG\_HEAD}, so this often works
straight away:

\begin{shell}
git reset --hard ORIG_HEAD
\end{shell}

\subsection{When the reflog is not enough}

If a commit is not in the reflog and not on any branch, it is still in the object
database until garbage collection runs. Find orphans with:

\begin{shell}
git fsck --lost-found
\end{shell}

This lists dangling commits and blobs. Inspect one with \co{git show <sha>}, and if
it is what you wanted, \co{git switch -c recovered <sha>}.

\begin{deeper}
Unreferenced objects survive for ninety days by default, controlled by
\co{gc.reflogExpire} and \co{gc.pruneExpire}. Git runs \co{gc} automatically when
loose objects pile up.

Which means the honest version of ``git never loses anything'' is: git does not
lose anything for about three months, as long as it was committed at least once.
Everything in this phase depends on that word \emph{committed}. Committing early
and often is not a discipline for other people's benefit, it is what makes you
recoverable.
\end{deeper}

\section{Stash}

Sometimes you do not want to undo work, you want to put it down for ten minutes.

\begin{shell}
git stash
git stash pop
\end{shell}

\co{stash} takes your uncommitted changes, saves them, and gives you a clean working
directory. \co{pop} puts them back and removes the saved copy.

Useful in practice:

\vspace{0.5em}
\begin{tabularx}{\linewidth}{@{}L{5.6cm} X@{}}
\toprule
\rhtablehead{Command} & \rhtablehead{What it does} \\
\midrule
\co{git stash push -m "half a fix"} & Save with a label you will recognise \\
\co{git stash list} & Show everything you have stashed \\
\co{git stash show -p stash@\{1\}} & Look at one as a diff before applying it \\
\co{git stash apply} & Restore but keep the stash as well \\
\co{git stash pop} & Restore and delete the stash \\
\co{git stash -u} & Include untracked files \\
\co{git stash branch fix} & Make a new branch from a stash \\
\co{git stash drop stash@\{0\}} & Delete one \\
\co{git stash clear} & Delete all of them \\
\bottomrule
\end{tabularx}
\vspace{0.5em}

\begin{trap}
The stash is a stack and it is easy to treat as a filing cabinet. Six unlabelled
stashes from three weeks ago are useless, because you will never work out which is
which and you will never dare delete any of them.

If work is worth keeping for more than an hour, it is worth a branch and a commit.
Branches have names.
\end{trap}

\begin{note}
Stashes are reachable through \co{git fsck} even after \co{git stash drop},
because they are ordinary commit objects. A dropped stash is recoverable for the
usual ninety days, though finding it takes effort.
\end{note}

\section{Choosing the right undo}

\vspace{0.5em}
\begin{tabularx}{\linewidth}{@{}L{6.4cm} X@{}}
\toprule
\rhtablehead{Situation} & \rhtablehead{Command} \\
\midrule
Edited a file, want the committed version & \co{git restore <file>} \\
Staged something by mistake & \co{git restore -{}-staged <file>} \\
Untracked junk to delete & \co{git clean -n} then \co{git clean -fd} \\
Bad commit message, not pushed & \co{git commit -{}-amend} \\
Forgot a file in the last commit & \co{git add <f>} then \co{git commit -{}-amend -{}-no-edit} \\
Committed too early, want the work back & \co{git reset -{}-soft HEAD\textasciitilde 1} \\
Want the last commit gone completely & \co{git reset -{}-hard HEAD\textasciitilde 1} \\
Bad commit already pushed to main & \co{git revert <sha>} \\
Deleted a branch you needed & \co{git reflog} then \co{git switch -c <name> <sha>} \\
Rebase went wrong & \co{git reset -{}-hard ORIG\_HEAD} \\
Need to set work aside briefly & \co{git stash} \\
Cannot find a commit anywhere & \co{git fsck -{}-lost-found} \\
\bottomrule
\end{tabularx}
\vspace{0.5em}

\section{What genuinely cannot be recovered}

Short list, and worth knowing exactly:

\begin{enumerate}
  \item \textbf{Changes never committed or stashed}, destroyed by
    \co{git restore}, \co{git checkout} or \co{git reset --hard}. Git never held a
    copy.
  \item \textbf{Untracked files} removed by \co{git clean}. Same reason.
  \item \textbf{Anything already garbage collected}, roughly ninety days after it
    became unreachable.
  \item \textbf{A repository you deleted from disk}, unless it was pushed somewhere.
\end{enumerate}

Every item on that list is a variation of the same sentence: git can only give you
back what it was given. The habit that follows is commit early, commit often, push
regularly. A commit costs nothing and buys you three months of insurance.

\checkyourself{
\begin{enumerate}
  \item Name the three trees, and say which ones \co{--soft}, \co{--mixed} and
    \co{--hard} each touch.
  \item You committed too early and want the changes back, staged. Which command?
  \item Why is \co{git revert} correct for a pushed commit when \co{git reset} is not?
  \item You deleted a branch an hour ago. Walk through recovering it.
  \item Which two commands in this phase can genuinely destroy work, and what do
    they have in common?
\end{enumerate}}

\vspace{1em}
{\color{rhborder}\rule{\linewidth}{0.8pt}}

\textbf{Next:} Phase 5, \emph{Rewriting history}, where you stop undoing commits and
start reshaping them: interactive rebase, squashing, splitting a commit in half, and
moving a commit from one branch to another.

\end{document}
